\documentclass[11pt,a4paper]{article}

\usepackage[utf8]{inputenc}
\usepackage[T1]{fontenc}
\usepackage{lmodern}
\usepackage[margin=2.5cm]{geometry}
\usepackage{amsmath,amssymb}
\usepackage{array}
\usepackage{booktabs}
\usepackage{tabularx}
\usepackage{caption}
\captionsetup[table]{position=top,skip=6pt,labelfont=bf,font=small}
\usepackage{enumitem}
\setlist[description]{style=unboxed}
\usepackage{microtype}
\usepackage[
  colorlinks=true,
  linkcolor=black,
  citecolor=black,
  urlcolor=black
]{hyperref}
\usepackage[capitalize,noabbrev]{cleveref}

\newcommand{\Cl}{\mathcal{C}\ell}
\newcommand{\R}{\mathbb{R}}
\newcommand{\C}{\mathbb{C}}
\newcommand{\Z}{\mathbb{Z}}

\title{\Large\bfseries Work Plan: Linear Spaces, Symmetries, and Clifford
  Algebras\\
  in Theoretical Physics}
\author{Clifford Seminar Group\\
  \small Author: Julian L. Avila-Martinez\\
  \small Universidad Distrital Francisco Jos\'{e} de Caldas}
\date{November 2026}

\begin{document}

\maketitle

\section{Conceptual Scope and Objective}
\label{sec:scope}

This work plan outlines a problem-driven study of the algebraic structures
underlying theoretical and mathematical physics. Rather than treating abstract
algebra in isolation, each mathematical formalism is introduced as the natural
tool required to solve and clarify a fundamental physical problem.

The progression begins with coordinate-free linear spaces and spectral
decomposition, applying them to decouple finite systems of coupled harmonic
oscillators before passing to infinite-dimensional function spaces to show how
the classical wave equation decomposes into an infinite family of independent
harmonic oscillators. It then examines how continuous symmetries of quadratic
linear spaces generate Lie algebras, and how completing a quadratic linear
space into a universal Clifford algebra $\Cl(V,q)$ embeds the Lie algebra
$\mathfrak{so}(p,q)$ directly within the subspace of bivectors. This machinery
is subsequently applied to unravel the hidden $\mathrm{SO}(4)$ dynamical
symmetry of the Kepler central force problem and the algebraic structure of
relativistic field equations in Minkowski spacetime.

\section{Thematic Modules and Problem Progression}
\label{sec:modules}

\subsection{Module I: The Linear Space Backbone and Oscillatory Systems}
\label{subsec:mod1}

\begin{description}
  \item[Topic 1: Linear Spaces, Dual Pairings, and Metric Structure.]
    Axiomatic formulation of linear spaces $V$ over $\R$ and $\C$ without
    reference to coordinate bases or matrices. Construction of the dual linear
    space $V^{*}$, canonical bilinear pairings, and the structural distinction
    between vectors and covectors in physical measurements. Introduction of
    symmetric bilinear forms, quadratic forms $q(v)$, the polarization
    identity, and metric signature $(p,q)$ as the geometric foundation of
    Euclidean space and relativistic spacetime.

  \item[Topic 2: Spectral Theory and Coupled Harmonic Oscillators.]
    Linear operators $T \in \mathcal{L}(V)$, invariant subspaces, and adjoint
    operators defined through the non-degenerate bilinear form. Proof of the
    spectral theorem for self-adjoint operators via invariant subspaces rather
    than determinants.
    \emph{Guiding Problem:} Formulation of $N$ coupled harmonic oscillators
    $\ddot{u} + Ku = 0$ as an operator equation on a finite-dimensional linear
    space, showing how orthogonal projections onto invariant eigenspaces
    decouple the dynamics into independent normal modes.

  \item[Topic 3: Infinite-Dimensional Spaces and the Wave Equation.]
    Extension of linear space concepts from discrete oscillator lattices to
    infinite-dimensional function spaces equipped with an integral inner
    product. The Fourier transformation as a unitary change of basis that
    diagonalizes translation-invariant differential operators.
    \emph{Guiding Problem:} Spectral decomposition of the classical wave
    equation $(\partial_{t}^{2} - c^{2}\nabla^{2})\Phi = 0$ and the
    Klein--Gordon equation, proving that in momentum space a continuous field
    is isomorphic to an infinite collection of decoupled harmonic oscillators
    $\ddot{\tilde{\Phi}}_{k} + \omega_{k}^{2}\tilde{\Phi}_{k} = 0$.
\end{description}

\subsection{Module II: Continuous Symmetries, Lie Algebras, and Clifford
  Algebras}
\label{subsec:mod2}

\begin{description}
  \item[Topic 4: Infinitesimal Transformations and Lie Algebras.]
    One-parameter continuous groups of transformations preserving the quadratic
    form $q(v)$ on a linear space. Generation of finite transformations via the
    operator exponential $\exp(\theta A)$ and the conservation requirement of
    skew-adjointness $A^{\dagger} = -A$. Emergence of the Lie bracket
    $[A,B] = AB - BA$ from the failure of transformations to commute, the
    Jacobi identity, and the derivation of the Lie algebras $\mathfrak{so}(3)$
    and $\mathfrak{su}(2)$ from rotational invariance.

  \item[Topic 5: The Universal Clifford Algebra of a Quadratic Space.]
    Motivation from the algebraic limitation of a bare linear space, where
    vectors cannot be multiplied directly to encode both projection and planar
    orientation. Construction of the universal Clifford algebra $\Cl(V,q)$ as
    the quotient of the tensor algebra $T(V)$ by the two-sided ideal
    $I_{q} = ( v \otimes v - q(v)1 )$. Decomposition of the
    geometric product $uv = u \cdot v + u \wedge v$, $\Z_{2}$-grading,
    multivector blade projections, grade involution $\hat{u}$, and reversion
    $\tilde{u}$.

  \item[Topic 6: Bivector Lie Algebras, Rotors, and Spin Groups.]
    Demonstration that the grade-2 subspace of bivectors $\Cl^{2}(V,q)$ closes
    under the commutator product $A \times B = \frac{1}{2}(AB - BA)$ and is
    isomorphic to the orthogonal Lie algebra $\mathfrak{so}(p,q)$.
    Exponentiation of simple bivectors into rotors
    $R = \exp(-B/2) \in \mathrm{Spin}(p,q)$, coordinate-free orthogonal
    transformations via the double-sided action $v \mapsto Rv\tilde{R}$, and
    the topological double cover $\mathrm{Spin}(3) \cong \mathrm{SU}(2) \to
    \mathrm{SO}(3)$ without matrix representations.
\end{description}

\subsection{Module III: Central Force Dynamics and the
  \texorpdfstring{$\mathrm{SO}(4)$}{SO(4)} Kepler Symmetry}
\label{subsec:mod3}

\begin{description}
  \item[Topic 7: Central Forces and the Angular Momentum Bivector.]
    Formulation of Newtonian mechanics in the Euclidean Clifford algebra
    $\Cl(3,0)$ for an isotropic central interaction
    $\dot{p} = -f((r^{2})^{1/2})r$. Replacement of the axial cross-product
    vector by the intrinsic orbital angular momentum bivector $L = r \wedge p$
    representing the oriented orbital plane. Coordinate-free proof of
    $\dot{L} = 0$, the exterior condition for planar motion $r \wedge L = 0$,
    and duality via the Euclidean pseudoscalar $I_{3}$.

  \item[Topic 8: The Kepler Problem and the Laplace--Runge--Lenz Invariant.]
    Specialization to the inverse-square interaction
    $\dot{p} = -k(r^{2})^{-3/2}r$. Direct geometric multiplication of linear
    momentum and the angular momentum bivector $pL = p \cdot L + p \wedge L$ to
    isolate the conserved Laplace--Runge--Lenz vector
    $A = \frac{1}{mk}pL - (r^{2})^{-1/2}r$. Derivation of $\dot{A} = 0$, the
    coplanarity relation $A \wedge L = 0$, and immediate algebraic extraction
    of conic sections from the inner contraction $r \cdot A$.

  \item[Topic 9: Hidden $\mathrm{SO}(4)$ Dynamical Symmetry.]
    Analysis of the algebraic brackets between the components of the angular
    momentum bivector $L$ and the rescaled Laplace--Runge--Lenz vector $A$ on
    bound-state energy manifolds. Decoupling of the conserved generators into
    two commuting $\mathfrak{su}(2)$ algebras, establishing the dynamical
    symmetry $\mathfrak{so}(4) \cong \mathfrak{su}(2) \oplus \mathfrak{su}(2)$
    responsible for closed trajectories and spectral degeneracy, alongside a
    comparison with the dynamical symmetry of the isotropic harmonic
    oscillator.
\end{description}

\subsection{Module IV: Relativistic Field Equations in Spacetime Algebra}
\label{subsec:mod4}

\begin{description}
  \item[Topic 10: Spacetime Algebra $\Cl(1,3)$ and Lorentz Symmetries.]
    Construction of the Clifford algebra of Minkowski spacetime with signature
    $(+,-,-,-)$ generated by $\{\gamma_{\mu}\}$. The observer-dependent
    spacetime split induced by a timelike vector $\gamma_{0}$, mapping
    spacetime bivectors into relative spatial vectors
    $\sigma_{k} = \gamma_{k}\gamma_{0}$ and spatial bivectors
    $I_{4}\sigma_{k}$. Unification of spatial rotations and relativistic boosts
    as rotor transformations $x \mapsto Rx\tilde{R}$ generated by the
    six-dimensional bivector Lie algebra
    $\Cl^{2}(1,3) \cong \mathfrak{so}(1,3)$.

  \item[Topic 11: Electrodynamics and Relativistic Field Oscillators.]
    Introduction of the coordinate-free spacetime vector derivative
    $\nabla = \gamma^{\mu}\partial_{\mu}$ satisfying
    $\nabla^{2} = \partial_{t}^{2} - \nabla_{\mathrm{spatial}}^{2}$. Synthesis
    of the electric vector field and magnetic bivector field into a single
    Faraday bivector $F = E + I_{4}B$, condensing the four Maxwell equations
    into the single multivector equation $\nabla F = J$. Application of the
    Module I Fourier mode expansion to the free-field equation $\nabla F = 0$,
    decomposing the electromagnetic field into relativistic harmonic
    oscillators.

  \item[Topic 12: Algebraic Foundations of the Relativistic Dirac Field.]
    First-order factorization of the relativistic Klein--Gordon wave operator
    $\nabla^{2}\Psi = -m^{2}\Psi$ using the geometric product of the spacetime
    derivative. Representation of relativistic spinor fields within the real
    even subalgebra $\Cl^{+}(1,3)$, demonstrating how complex phase
    transformations correspond geometrically to rotations generated by the
    spin-plane bivector $\gamma_{2}\gamma_{1}$.
\end{description}

\section{Summary of the Problem-Driven Progression}
\label{sec:summary}

A synthesis of the twelve topics, pairing each mathematical structure with its
corresponding physical problem, is presented in \cref{tab:summary}.

\begin{table}[htbp]
  \centering
  \caption{Synthesis of mathematical structures and guiding physical problems.}
  \label{tab:summary}
  \small
  \renewcommand{\arraystretch}{1.25}
  \begin{tabularx}{\textwidth}{%
    @{}%
    c%
    c%
    >{\raggedright\arraybackslash}p{4.8cm}%
    >{\raggedright\arraybackslash}X%
    @{}%
  }
    \toprule
    Topic & Module & Mathematical Structure & Guiding Physical Problem \\
    \midrule
    1  & I   & Linear spaces, duals, quadratic forms &
      Metric geometry of physical states and observables \\
    2  & I   & Self-adjoint operators, spectral theorem &
      Decoupling $N$ coupled harmonic oscillators into normal modes \\
    3  & I   & Function spaces, Fourier diagonalization &
      Wave equation $(\partial_{t}^{2} - c^{2}\nabla^{2})\Phi = 0$ as infinite
      harmonic oscillators \\
    \midrule
    4  & II  & Exponential map, commutators, Lie algebras &
      Continuous rotational invariance and $\mathfrak{so}(3) \cong
      \mathfrak{su}(2)$ \\
    5  & II  & Quotient tensor algebra $\Cl(V,q)$ &
      Unifying metric projection and planar orientation in products of
      vectors \\
    6  & II  & Bivector algebra $\Cl^{2}(V,q)$, $\mathrm{Spin}(p,q)$ &
      Coordinate-free spatial rotations $v \mapsto Rv\tilde{R}$ and spinor
      double cover \\
    \midrule
    7  & III & Exterior product and bivector conservation &
      Central force dynamics $\dot{p} = -f((r^{2})^{1/2})r$ and planar
      confinement $L = r \wedge p$ \\
    8  & III & Geometric product $pL = p \cdot L + p \wedge L$ &
      Inverse-square Kepler problem and Laplace--Runge--Lenz invariant $A$ \\
    9  & III & Lie bracket closure on energy manifolds &
      Hidden $\mathfrak{so}(4) \cong \mathfrak{su}(2) \oplus \mathfrak{su}(2)$
      symmetry and harmonic oscillator duality \\
    \midrule
    10 & IV  & Minkowski Clifford algebra $\Cl(1,3)$ &
      Lorentz boosts, spatial rotations, and observer spacetime split \\
    11 & IV  & Vector derivative $\nabla$ and Faraday bivector $F$ &
      Unification of Maxwell equations $\nabla F = J$ and radiation modes \\
    12 & IV  & Even subalgebra $\Cl^{+}(1,3)$ &
      Factorization of the relativistic wave operator and the Dirac field \\
    \bottomrule
  \end{tabularx}
\end{table}

\end{document}
